\documentclass[12pt]{article}
\usepackage{amsmath, amssymb}
\usepackage{geometry}
\usepackage{tcolorbox}
\usepackage{enumitem}
\usepackage{hyperref}

\geometry{a4paper, margin=1in}
\setlength{\parskip}{1em}
\setlength{\parindent}{0em}

\title{Metric Spaces and Examples}
\author{}
\date{}

\begin{document}

\maketitle

수업에서 다루는 dataset은 앞으로 ""finite metric measure space""라 생각하자.

\section*{Metric Space}

\subsection*{Definition}

\begin{tcolorbox}[colback=gray!5!white, colframe=gray!75!black, title=Definition]
A metric space $(X, d_X)$ is a pair consisting of a set $X$ and a function $d_X : X \times X \to \mathbb{R}_+$ such that
\[
\begin{aligned}
    & d_X(x, y) = 0 \iff x = y \\
    & d_X(x, y) = d_X(y, x) \\
    & d_X(x, y) + d_X(y, z) \ge d_X(x, z)
\end{aligned}
\]
A function $d_X$ is referred to as a ``metric'' or ``distance''.
\end{tcolorbox}

$d_X(x, y) + d_X(y, z) \ge d_X(x, z)$는 ``최단거리''에 대한 일반화이다.

\subsection*{Example}

\begin{tcolorbox}[colback=gray!5!white, colframe=gray!75!black, title=Example]
The following are some metrics on $\mathbb{R}^n$.

\textbf{1.} $d_1(x, y) = \Vert x - y \Vert_p := \left( \sum_{i=1}^n \vert x_i - y_i \vert \right)^{1/p}$ for $p \in \{1, 2, \cdots\}$.
\end{tcolorbox}

\end{document}
